\documentclass[11pt,a4paper]{article}
\newcommand{\sheetnumber}{1}
\newcommand{\submissionauthor}{Tobias Steindl}
% Shared preamble for exercise submissions. Load after \documentclass.
\usepackage[T1]{fontenc}
\usepackage[utf8]{inputenc}
\usepackage{lmodern}
\usepackage[a4paper,margin=25mm]{geometry}
\usepackage{amsmath,amssymb}
\usepackage{fancyhdr}
\usepackage{hyperref}
\hypersetup{hidelinks}
\setlength{\parindent}{0pt}
\setlength{\parskip}{0.5em}
\setlength{\headheight}{14pt}
\pagestyle{fancy}
\fancyhf{}
\fancyhead[L]{Numerical Methods for Elliptic Problems}
\fancyhead[R]{Sheet \sheetnumber}
\fancyfoot[C]{\thepage}
\newcommand{\R}{\mathbb{R}}
\newcommand{\dx}{\,\mathrm{d}x}
\newcommand{\ds}{\,\mathrm{d}s}
\newcommand{\Div}{\operatorname{div}}
\newcommand{\submissionheader}{%
  \begin{center}
    {\Large\bfseries Numerical Methods for Elliptic Problems\par}
    {\large Exercise Sheet \sheetnumber\par}
    Johannes Kepler University Linz --- Winter Semester 2026/27\par
    \ifx\submissionauthor\empty\else\submissionauthor\par\fi
  \end{center}\medskip}

\usepackage{graphicx}
\begin{document}
\begingroup
\setlength{\parskip}{0pt}
\submissionheader
\endgroup
\enlargethispage{\baselineskip}

\section*{Exercise 1 (Vector identities, integration-by-parts)}
\textbf{(a)}
\[
\Div(\vec b u)=\sum_{i=1}^d\partial_i(ub_i)
\overset{\mathrm{PR}}{=}\sum_{i=1}^d\bigl((\partial_i u)b_i+u\partial_i b_i\bigr)
=\vec b\cdot\nabla u+u\Div\vec b.
\]
\textbf{(b)} Gau\ss:
\[
\int_{\partial\Omega}u\vec b\cdot\vec n\ds
=\int_\Omega\Div(u\vec b)\dx
\overset{(a)}{=}\int_\Omega\vec b\cdot\nabla u\dx+\int_\Omega u\Div\vec b\dx.
\]
\textbf{(c)} In (b) choose $\vec b=\nabla v$, then
\[
\int_\Omega\nabla u\cdot\nabla v\dx
=-\int_\Omega u\underbrace{\Div(\nabla v)}_{\Delta v}\dx
+\int_{\partial\Omega}u\underbrace{\nabla v\cdot\vec n}_{\partial_n v}\ds.
\]
Exchange $u,v$ and use symmetry of inner product:
\[
\int_\Omega\nabla u\cdot\nabla v\dx
=-\int_\Omega\Delta u\,v\dx+\int_{\partial\Omega}\partial_n u\,v\ds.
\]

\section*{Exercise 2 (Fundamental lemma 1)}
\textbf{1. $\Leftarrow$} Clear.

\textbf{$\Rightarrow$} Assume for the sake of contradiction there exists
$x_0\in\Omega$ with $f(x_0)\ne0$. Without loss $f(x_0)>0$.
Due to the continuity of $f$, there exists $\delta>0$ such that $f(x)>0$
for $x\in B_\delta(x_0)\subset\Omega$ with $\delta$ sufficiently small
($\Omega$ is open!). Now choose $r<\delta$ such that
$\overline B_r(x_0)\subset\Omega$. It holds:
\[
0=\int_{\overline B_r(x_0)}f(x)\dx
\ge\int_{B_\delta(x_0)}\underbrace{f(x)}_{>0}\dx>0.\qquad\bot
\]
\textbf{2. $\Leftarrow$}
$f=0$ a.e. $\Rightarrow\int_\omega f\dx=0$ (clear).

\textbf{$\Rightarrow$} We show that $\int_N f(x)\dx=0$ for all null sets $N$.
For this choose a ``subset on which $f$ is strictly positive'':
\[
E_{n,R}=\left\{x\in\Omega\cap B_R(0):f(x)\ge\frac1n\right\}.
\]
Clearly for $x\in E_{n,R}$: $f(x)\ge\frac1n>0$, for all $n,R$ by definition.
Also $E_{n,R}$ measurable (clear) and bounded due to
``$\cap B_R(0)$'' for $R<\infty$. Now,
\[
0\overset{\mathrm{Ass.}}{=}\int_{E_{n,R}}\underbrace{f(x)}_{\ge1/n}\dx
\ge\underbrace{\frac1n}_{>0}\underbrace{|E_{n,R}|}_{\ge0}
\quad\Rightarrow\quad |E_{n,R}|=0.
\]
Analogously, for
$E'_{n,R}=\{x\in\Omega\cap B_R(0):f(x)\le-\frac1n\}$
it follows that $f(x)<\frac1n<0$, $\mu(E'_{n,R})=0$. Now,
\[
A_{n,R}=E_{n,R}\cup E'_{n,R}
\]
and
\[
B=\bigcup_{n,R\in\mathbb N}A_{n,R},\qquad\mu(B)=0\quad\text{(countable union)}.
\]
Equivalently, $f=0$ a.e.

\newpage
\section*{Exercise 3 (Fundamental lemma, 2nd version)}
\textbf{$\Leftarrow$}
\[
f(x)=0\Rightarrow f(x)v(x)=0\quad\forall v\in C(\overline\Omega)
\Rightarrow 0=\int_\Omega fv\dx\quad\forall v.
\]
\textbf{$\Rightarrow$} Suppose $f(x_0)\ne0$ for $x_0\in\Omega$.
W.l.o.g. $f(x_0)>0$. Continuity:
$\overline B_r(x_0)\subset\Omega$ with $f(x)>0$ for all
$x\in\overline B_r(x_0)$. Define
\[
v(x)=\max\left\{1-\frac{|x-x_0|}{r},0\right\},\qquad x\in\overline\Omega.
\]
Then $v\ge0$, $v(x_0)=1$, $v(x)=0$ outside $B_r(x_0)$ and $v$ is continuous.
In particular $v|_{\partial\Omega}=0$, so $v$ is admissible. Furthermore
\[
f(x)v(x)>0\quad\text{on }B_r(x_0),\qquad
f(x)v(x)=0\quad\text{outside }B_r(x_0).
\]
Thus,
\[
\int_\Omega f(x)v(x)\dx=\int_{B_r(x_0)}\underbrace{f(x)v(x)}_{>0}\dx>0.\qquad\bot
\]
\textbf{$L^2$ case: $\Leftarrow$}
Clearly $f=0$ a.e. $\Rightarrow\int_\Omega fv\dx=0$ for all $v\in L^2(\Omega)$.
As $\Omega$ is bounded, if $f\in C_c^\infty(\Omega)$ then $f\in L^2(\Omega)$.
Thus $C_c^\infty(\Omega)\subset L^2(\Omega)$ and thus the leftmost
implication ``$\Leftarrow$'' holds.

\textbf{1. $\Rightarrow$} Suppose $\int_\Omega f\varphi\dx=0$
for all $\varphi\in C_c^\infty(\Omega)$.
Let $v\in L^2(\Omega)$ arbitrary but fixed. We show $\int_\Omega fv\dx=0$.
Due to density of $C_c^\infty(\Omega)$ in $L^2(\Omega)$ there exists a
sequence $(\varphi_n)\subset C_c^\infty$:
\[
\|\varphi_n-v\|_{L^2(\Omega)}\to0.
\]
Due to C--S-inequality:
\[
|\langle f,v-\varphi_n\rangle|_{L^2}
\le\|f\|_{L^2}\underbrace{\|v-\varphi_n\|_{L^2}}_{\to0}\to0.
\]
Due to assumption: $\langle f,\varphi_n\rangle=0$ for all $n$,
as $\varphi_n\in C_c^\infty$.
\begin{align*}
\langle f,v\rangle
&=\langle f,v-\varphi_n+\varphi_n\rangle\\
&=\underbrace{\langle f,v-\varphi_n\rangle}_{\le |\cdot|}
+\underbrace{\langle f,\varphi_n\rangle}_{0}
=\langle f,v-\varphi_n\rangle\to0.
\end{align*}
Hence, it must also hold that $\langle f,v\rangle=0$.

\textbf{2. $\Rightarrow$} Suppose $\int_\Omega fv\dx=0$ for all $v\in L^2(\Omega)$.
As $f\in L^2$, in particular $v=f$:
\[
0=\int_\Omega f^2\dx=\|f\|_{L^2}^2\quad\Rightarrow\quad f=0.
\]

\newpage
\section*{Exercise 4}
\textbf{(a)} We compute the first variation
\begin{align*}
\delta J(u)[v]
&=\lim_{t\to0}\frac1t\bigl[J(u+tv)-J(u)\bigr]\\
&=\lim_{t\to0}\frac1{2t}\Bigl[
\int_\Omega a|\nabla(u+tv)|^2\dx
-\int_\Omega fu\dx-t\int_\Omega fv\dx\\
&\hspace{38mm}-\int_\Omega a|\nabla u|^2\dx+\int_\Omega fu\dx\Bigr].
\end{align*}
\[
\nabla(u+tv)=\nabla u+t\nabla v,\qquad
|\nabla u+t\nabla v|^2
=|\nabla u|^2+2t\nabla u\cdot\nabla v+t^2|\nabla v|^2.
\]
\begin{align*}
\delta J(u)[v]
&=\lim_{t\to0}\frac1{2t}\Bigl[
\int_\Omega a|\nabla u|^2\dx+2t\int_\Omega a\nabla u\cdot\nabla v\dx
+t^2\int_\Omega a|\nabla v|^2\dx\\
&\hspace{38mm}-t\int_\Omega fv\dx-\int_\Omega a|\nabla u|^2\dx\Bigr]\\
&=\int_\Omega a\nabla u\cdot\nabla v\dx-\int_\Omega fv\dx
+\underbrace{\lim_{t\to0}t\int_\Omega a|\nabla v|^2\dx}_{0}.
\end{align*}
\[
\int_\Omega a|\nabla v|^2\dx
\le\|a\|_{L^\infty(\Omega)}\|\nabla v\|_{L^2}^2,
\qquad\text{because we can choose }v\in H_0^1(\Omega).
\]
Setting the first variation to zero we obtain:
\[
\int_\Omega a\nabla u\cdot\nabla v\dx=\int_\Omega fv\dx
\qquad\forall v\in H_0^1.
\]
Comparison to min surface:
\[
\int_\Omega\frac1{\sqrt{1+|\nabla u|^2}}\nabla u\cdot\nabla v\dx=0,
\qquad\forall v:\ v|_{\partial\Omega}=0.
\]
Here, we are linear in $\nabla u$, $\nabla v$: $a\nabla u\cdot\nabla v$.
MS is non-linear: $\frac1{\sqrt{1+|\nabla u|^2}}\nabla u\cdot\nabla v$.
However, we have a RHS $\int_\Omega fv\dx$.

\textbf{(b)} Using (a) from Ex 1 and restricting to
$v\in C_c^\infty\subset H_0^1$:
\[
\Div\bigl((a\nabla u)v\bigr)
=(a\nabla u)\cdot\nabla v+v\Div(a\nabla u),\qquad\forall v\in C_c^\infty.
\]
\[
\int_\Omega\Div\bigl((a\nabla u)v\bigr)\dx
=\underbrace{\int_\Omega a\nabla u\cdot\nabla v\dx}_{\int_\Omega fv\dx\ \text{(Teil a)}}
+\int_\Omega v\Div(a\nabla u)\dx.
\]
Gau\ss:
\[
\int_{\partial\Omega}\bigl((a\nabla u)v\bigr)\cdot\vec n\ds=0
\qquad\text{($v=0$ due to $v\in H_0^1$)}.
\]
\begin{align*}
0&=\int_\Omega fv\dx+\int_\Omega v\Div(a\nabla u)\dx\\
&=\int_\Omega v\bigl(f+\Div(a\nabla u)\bigr)\dx
\qquad\forall v\in C_c^\infty.
\end{align*}
By 3 we obtain:
\[
f+\Div(a\nabla u)=0\quad\Leftrightarrow\quad-\Div(a\nabla u)=f
\quad\text{in }\Omega,
\]
and boundary condition
\[
u=0\quad\text{on }\partial\Omega,
\]
from the definition of the minimization problem.

\section*{Exercise 5}
Let $u\in C^2(\Omega)\cap C(\overline\Omega)$ be a classical solution of
\[
-\Delta u=0\quad\text{in }\Omega,\qquad u=g\quad\text{on }\partial\Omega.
\]
Choose $w=u-u_g$, then it is admissible:
\[
-\Delta w=\underbrace{-\Delta u}_{0}+\Delta u_g=\Delta u_g=f
\quad\text{in }\Omega.
\]
\[
w=\underbrace{u}_{g}-\underbrace{u_g}_{g}=0.
\]
Furthermore,
\[
u=u_g+w=u_g+u-u_g=u.
\]

\section*{Exercise 7}
\begin{center}
\includegraphics[page=12,trim=72 250 235 204,clip,width=0.65\textwidth]{Ex 1.pdf}
\end{center}
\end{document}
